\documentclass[10pt, letterpaper]{article}

% Packages:
\usepackage[
    ignoreheadfoot,
    top=1.6 cm,
    bottom=1.6 cm,
    left=2 cm,
    right=2 cm,
    footskip=1.0 cm,
]{geometry}
\usepackage[explicit]{titlesec}
\usepackage{tabularx}
\usepackage{array}
\usepackage[dvipsnames]{xcolor}
\definecolor{primaryColor}{RGB}{0, 79, 144}
\usepackage{enumitem}
\usepackage{fontawesome5}
\usepackage{amsmath}
\usepackage[
    pdftitle={Currículum Vitae},
    pdfauthor={Jesús Antonio Romero Godoy},
    colorlinks=true,
    urlcolor=primaryColor
]{hyperref}
\usepackage[pscoord]{eso-pic}
\usepackage{calc}
\usepackage{bookmark}
\usepackage{lastpage}
\usepackage{changepage}
\usepackage{paracol}
\usepackage{ifthen}
\usepackage{needspace}
\usepackage{iftex}
\usepackage{etoolbox}

\ifPDFTeX
    \input{glyphtounicode}
    \pdfgentounicode=1
    \usepackage[T1]{fontenc}
    \usepackage[utf8]{inputenc}
    \usepackage{lmodern}
\fi

\usepackage[default, type1]{sourcesanspro}

\AtBeginEnvironment{adjustwidth}{\partopsep0pt}
\pagestyle{empty}
\setcounter{secnumdepth}{0}
\setlength{\parindent}{0pt}
\setlength{\topskip}{0pt}
\setlength{\columnsep}{0.15cm}
\makeatletter
\let\ps@customFooterStyle\ps@plain
\patchcmd{\ps@customFooterStyle}{\thepage}{
    \color{gray}\textit{\small Jesús Antonio Romero Godoy}
}{}{}
\makeatother
\pagestyle{customFooterStyle}

\titleformat{\section}{
    \needspace{4\baselineskip}
    \Large\color{primaryColor}
}{ }{ }{\textbf{#1}\hspace{0.15cm}\titlerule[0.8pt]\hspace{-0.1cm}}[]

\titlespacing{\section}{-1pt}{0.3 cm}{0.2 cm}

\newenvironment{highlights}{
    \begin{itemize}[
        label={-},
        topsep=0pt,
        parsep=0pt,
        partopsep=0pt,
        itemsep=0pt,
        leftmargin=0.35 cm,
        labelsep=0.15 cm
    ]
}{
    \end{itemize}
}

\newenvironment{onecolentry}{
    \begin{adjustwidth}{0.2 cm+0.00001 cm}{0.2 cm+0.00001 cm}
}{
    \end{adjustwidth}
}

\newenvironment{twocolentry}[2][]{
    \onecolentry
    \def\secondColumn{#2}
    \setcolumnwidth{\fill, 4.5 cm}
    \begin{paracol}{2}
}{
    \switchcolumn \raggedleft \secondColumn
    \end{paracol}
    \endonecolentry
}

\newenvironment{header}{
    \setlength{\topsep}{0pt}\par\kern\topsep\centering\color{primaryColor}\linespread{1.5}
}{
    \par\kern\topsep
}

\newcommand{\placelastupdatedtext}{
  \AddToShipoutPictureFG*{
    \put(
        \LenToUnit{\paperwidth-2 cm-0.2 cm+0.05cm},
        \LenToUnit{\paperheight-1.0 cm}
    ){\vtop{{\null}\makebox[0pt][c]{
        \small\color{gray}\textit{Julio 2026}
    }}}%
  }%
}

\begin{document}
    \placelastupdatedtext

%====================
% Header centrado
%====================
\begin{header}
    \fontsize{30 pt}{30 pt}\selectfont
    \textbf{Jesús Antonio Romero Godoy}

    \vspace{0.3 cm}

    \normalsize
    \faMapMarker* \hspace{0.13cm} Ciudad de México, México
    \kern 0.5 cm
    \faEnvelope[regular] \hspace{0.13cm} \href{mailto:jromerog01@outlook.com}{jromerog01@outlook.com}
    \kern 0.5 cm
    \faPhone* \hspace{0.13cm} +52 55 1652 2741
    \\
    \faGithub \hspace{0.13cm} \href{https://github.com/jromerog01}{jromerog01}
    \kern 0.5 cm
    \faLinkedin \hspace{0.13cm} \href{https://www.linkedin.com/in/jesus-antonio-romero-godoy-8a9032383/}{Jesus Romero}
    \kern 0.5 cm
    \faGlobe \hspace{0.13cm} \href{https://jromerog.dev/}{jromerog.dev}
\end{header}

\vspace{0.5 cm}

%====================
% Perfil Profesional
%====================
\section{Perfil Profesional}
\begin{onecolentry}
Estudiante de 7º semestre de Ciencias de la Computación en la UNAM, enfocado en desarrollo backend con Java y Spring Boot. Más de tres años programando, con experiencia en PostgreSQL, JPA/Hibernate, Docker y despliegue de servicios en producción en un servidor propio. Comprometido con la calidad del código, el aprendizaje continuo y el trabajo en equipo.
\end{onecolentry}

%====================
% Formación Académica
%====================
\section{Formación Académica}
\begin{twocolentry}{
    UNAM, Facultad de Ciencias\\
    Ago 2023 – May 2027
}
    \textbf{Licenciatura en Ciencias de la Computación} \\
    Promedio: 9.58/10 \\
    Avance en Créditos: 71\%
\end{twocolentry}

%====================
% Proyectos
%====================
\section{Proyectos}

\begin{twocolentry}{
    Java, Spring Boot,\\ PostgreSQL \& Docker
}
    \textbf{\href{https://github.com/jromerog01/StockFlow-APIRest}{StockFlow API — API REST de Inventarios y Punto de Venta}}
    \begin{highlights}
        \item Gestión de proveedores, productos y stock con punto de venta e historial de movimientos, persistencia con JPA/Hibernate sobre PostgreSQL.
        \item Documentada con Swagger/OpenAPI, contenerizada con Docker y en producción self-hosted en \mbox{\href{https://stockflow-api.jromerog.dev/swagger-ui/index.html}{stockflow-api.jromerog.dev}}.
    \end{highlights}
\end{twocolentry}

\vspace{0.15 cm}

\begin{twocolentry}{
    Ubuntu Server,\\ \& Docker 
}
    \textbf{\href{https://jromerog.dev/}{Servidor Personal / Home Lab}}
    \begin{highlights}
        \item Equipo con Ubuntu Server 24/7 que aloja mis proyectos con Docker y Docker Compose, administrando imágenes, volúmenes y redes de contenedores.
        \item Publicación mediante Cloudflare Tunnel (HTTPS en subdominios propios), acceso remoto seguro mediante SSH y Tailscale sin exponer puertos.
    \end{highlights}
\end{twocolentry}

\vspace{0.15 cm}

\begin{twocolentry}{
    PostgreSQL, PL/pgSQL\\ \& Spring Boot
}
    \textbf{\href{https://github.com/jromerog01/Proyecto-Final-Bases-de-Datos}{Sistema de Gestión de Eventos Multirrol}}
    \begin{highlights}
        \item Modelado relacional desde cero, con 34 tablas y 45 llaves foráneas, con disparadores y procedimientos almacenados en PL/pgSQL.
        \item API REST en Spring Boot con JDBC y panel web para el CRUD. Trabajo en equipo.
    \end{highlights}
\end{twocolentry}

\vspace{0.15 cm}

\begin{twocolentry}{
    Python, Django\\ \& Docker
}
    \textbf{\href{https://github.com/jromerog01/Sistema-Reservaciones-Luciernagas}{Sistema de Reservaciones de Hospedaje}}
    \begin{highlights}
        \item App Django de reservaciones multisede con mapa interactivo (Leaflet/OpenStreetMap), reglas de negocio de temporada y roles de Visitante, Cliente y Administrador.
        \item Desarrollo en equipo, con arquitectura basada en patrones Facade, Template, \mbox{Observer} y Adapter, autenticación con Google y despliegue contenerizado en \mbox{\href{https://lumina-fest.jromerog.dev/}{lumina-fest.jromerog.dev}}.
    \end{highlights}
\end{twocolentry}



%====================
% Habilidades Técnicas
%====================
\section{Habilidades Técnicas}
\begin{onecolentry}
    \textbf{Lenguajes:} Java 17/21, Python\\
    \textbf{Backend:} Spring Boot (JPA/Hibernate, Spring Web, Swagger/OpenAPI), Django\\
    \textbf{Bases de Datos:} PostgreSQL (modelo relacional, consultas avanzadas, PL/pgSQL), SQLite\\
    \textbf{DevOps / Infraestructura:} Docker, Docker Compose, Nginx, Cloudflare Tunnel, Tailscale, SSH\\
    \textbf{Sistemas Operativos:} Fedora Linux, Ubuntu Server, Windows\\
    \textbf{Herramientas:} Git, GitHub, Postman \\
    \textbf{Patrones de Diseño:} MVC, Factory, Builder, Strategy, Observer, Adapter

\end{onecolentry}

%====================
% Idiomas y Competencias Personales
%====================
\section{Idiomas y Competencias Personales}
\begin{onecolentry}
    \textbf{Idiomas:} Español (Nativo), Inglés – B2 (Upper Intermediate)\\
    \textbf{Competencias:} Trabajo en equipo, comunicación efectiva, adaptabilidad, atención al detalle y pensamiento analítico.
\end{onecolentry}



\end{document}
